\chapter{Forward Dynamics --- Propagation Methods}
\label{chap:7}
The forward dynamics problem presents us with two sets of unknowns: the joint accelerations and the joint constraint forces. It is usually not possible to solve for any of these unknowns locally at any one body; but it is possible to formulate equations that they must satisfy. Propagation methods work by calculating the coefficients of such equations locally, and propagating them to neighbouring bodies, until we eventually reach a point where we can solve the dynamics locally at one particular body. Having the solution at this one body, it becomes possible to solve the dynamics at neighbouring bodies, and so on, until the whole problem is solved. In principle, the process is akin to solving a linear equation by first triangularizing the coefficient matrix, and then solving it by back-substitution. Indeed, some propagation algorithms have been formulated via this analogy (e.g. see Ascher et al., 1997; Rosenthal, 1990; Saha, 1997). Propagation algorithms are more complicated than inertia-matrix algorithms, but they have a computational complexity of O(NB), which is the theoretical minimum for solving the forward dynamics problem. The articulated-body algorithm is the prime example of a propagation algorithm, and it is the fastest for calculating the forward dynamics of a kinematic tree. Most of this chapter is devoted to a description of this algorithm and its attendant concepts, like articulated-body inertia. The final two sections take a broader view of the subject, and present some of the techniques and equations that form the basis of other algorithms. 7.1 Articulated-Body Inertia Articulated-body inertia is the inertia that a body appears to have when it is part of a rigid-body system. Consider the rigid body B shown in diagram (a) of Figure 7.1. The relationship between the applied force, f, and the resulting

B B B’ f a f a f = Ia + p f = IAa + pA (a) (b) Figure 7.1: Comparing a rigid body (a) with an articulated body (b) acceleration, a, is given by the body’s equation of motion: f = Ia + p . (7.1) The coefficients I and p in this equation are the rigid-body inertia and bias force, respectively, of body B. Now consider the two-body system shown in diagram (b), in which a second body, B′, has been connected to B via a joint. The effect of this second body is to alter the acceleration response of B, so that the relationship between f and a is now given by f = IA a + pA . (7.2) This is an articulated-body equation of motion. It has the same algebraic form as Eq. 7.1, but different coefficients. The quantities IA and pA in this equation are the articulated-body inertia and bias force of body B in the two-body system comprising B, B′ and the connecting joint. An equation like 7.2 describes the acceleration response of a single rigid body, taking into account the dynamic effects of a specific set of other bodies and joints. We call this body a handle, and we call the system containing it an articulated body. Thus, the handle is the body to which I A and pA refer; and the articulated body is the rigid-body system, or subsystem, whose dynamic effects are included in IA and pA. In Figure 7.1(b), the handle is body B and the articulated body is the whole two-body system. Basic Properties Articulated-body inertias share the following properties with rigid-body inertias: they are symmetric, positive-definite matrices; they are mappings from M6 to F6; and they obey the same coordinate transformation rule. However, articulated-body inertias are mappings from acceleration to force, not velocity to momentum, so expressions like IAv and 1 2vTIAv, where v denotes a velocity, do not make sense. An articulated body may consist of just a single rigid body, in which case the articulated-body and rigid-body inertias and bias forces are the same. This provides a starting point for recursive calculation methods.

A1 A2 A B1 B2 IA = I1 A + I2 A pA = p1 A + p2 A I1 A, p1 A I2 A, p2 A IA, pA B Figure 7.2: Addition of articulated-body inertias and bias forces Articulated-body inertias depend only on the inertias of the member bodies and the constraints imposed by the member joints. Thus, they are functions of the joint position variables, but not the velocity variables or the various force terms. In this respect, they resemble generalized inertia matrices. The velocity terms, external force terms (other than f), and so on, appear only in the bias forces. In 3D space, 21 parameters are required to define a general articulatedbody inertia, compared with only 10 for a rigid-body inertia. In 2D space, 6 parameters are required to define a general articulated-body inertia, compared with only 4 for a rigid-body inertia. Note that 21 and 6 are the numbers of independent values in a symmetric 6 × 6 and 3 × 3 matrix, respectively. In general, the apparent mass of an articulated-body inertia varies with the direction of applied force, and it does not have a centre of apparent mass. (See Example 7.1.) Addition Addition is defined on articulated-body inertias, and the physical interpretation of this operation is illustrated in Figure 7.2. This figure shows two articulated bodies, A1 and A2, that are being joined to form a single articulated body, A. The connection is made by gluing bodies B1 and B2 together to form a new body, B. If the symbols IA i , pA i , IA and pA denote the articulated-body inertias and bias forces for Bi in Ai and B in A, respectively, then IA = IA 1 + IA 2 and pA = pA 1 + pA 2 . Inverse Inertia Equation 7.2 assumes the handle has six degrees of freedom. If this is not the case, then the more general equation a = ΦAf + bA (7.3)

must be used instead. In this equation, ΦA and bA are the handle’s articulatedbody inverse inertia and bias acceleration, respectively (cf. Eq. 2.72). An equation in this form always exists. If the handle has six degrees of freedom, then ΦA = (IA)−1 and bA = −ΦA pA. Otherwise, ΦA is singular, and its rank equals the handle’s degree of motion freedom. Articulated-body inverse inertias are symmetric, positive-semidefinite matrices, and they transform like rigid-body inverse inertias. Multiple Handles It is possible for an articulated body to have more than one handle. If we number the handles from 1 to h, and define the vectors fi and ai to be the force and acceleration associated with handle i, then Eq. 7.3 generalizes to

a1 a2 ... ah

=

ΦA 1 ΦA 12 · · · ΦA 1h ΦA 21 ΦA 2 · · · ΦA 2h ... ... ... ... ΦA h1 ΦA h2 · · · ΦA h

f1 f2 ... fh

+

bA 1 bA 2 ... bA h

. (7.4) In this equation, ΦA i and bA i are the articulated-body inverse inertia and bias acceleration for handle i, and ΦA ij is the cross-coupling term that describes the acceleration caused at handle i by the application of a force to handle j. The coefficient matrix is symmetric and positive-semidefinite, and its rank equals the total number of independent degrees of motion freedom possessed by the set of handles. If the rank is 6h then Eq. 7.4 can be inverted to obtain a multi-handle equivalent of Eq. 7.2. O z y x Figure 7.3: An articulated body Example 7.1 Figure 7.3 shows an articulated body consisting of a box with a slot, and a cylinder that fits inside the slot. The cylinder has two degrees of freedom relative to the box: it can rotate about the x axis, and it can translate in the y direction. The contact between the two bodies is assumed to

be frictionless. Both bodies have their centres of mass at the origin. The box has a mass of m1 and a rotational inertia about its centre of mass of I1 (i.e., its rotational inertia matrix is I1 times the identity matrix); and the cylinder has a mass of m2 and a rotational inertia of I2. Let us treat the box as the handle, and calculate its articulated-body inertia. If we apply a force to the box along either the x or the z axis, then the two bodies will move as one, and the resulting acceleration will be 1/(m1 + m2) times the magnitude of the applied force; but if we apply a force along the y axis, then only the box will move, and the resulting acceleration will be 1/m1 times the magnitude of the applied force. Likewise, if a pure couple is applied to the box about either the y or the z axes, then the two bodies will rotate as one, and the apparent inertia will be I1 + I2; but a pure couple about the x axis will cause only the box to rotate, and the apparent inertia will be I1. The complete articulated-body inertia matrix is therefore IA =

I1 0 0 0 0 0 0 I1 + I2 0 0 0 0 0 0 I1 + I2 0 0 0 0 0 0 m1 + m2 0 0 0 0 0 0 m1 0 0 0 0 0 0 m1 + m2

. Observe that the apparent mass varies with direction: the box appears to have more mass in the x and z directions than in the y direction. In this particular example, the handle exhibits a centre of apparent mass, which is located at the origin. Any pure force applied to the handle along a line passing through this point causes a pure linear acceleration of the handle. Such a point does not exist in general. 7.2 Calculating Articulated-Body Inertias There are two basic strategies for calculating articulated-body inertias: projection and assembly. In the projection method, the equation of motion for the whole articulated body is projected onto the motion space of the handle (or handles). In the assembly method, an articulated body is assembled step-bystep from its component parts, and a sequence of articulated-body inertias are computed along the way. The projection method is commonly used to calculate operational-space inertias, which are used in robot control systems (Khatib, 1987, 1995); and the assembly method is used in O(n) dynamics algorithms. 7.2.1 The Projection Method An articulated body is a rigid-body system. It therefore has an equation of motion that can be expressed in generalized coordinates as H¨q + C = τ , (7.5)

where C accounts for every force acting on the system except the force acting on the handle. Let v, a and f denote the spatial velocity and acceleration of the handle, and the spatial force acting on it. If J is the Jacobian that maps ˙q to v according to v = J ˙q, then it defines the following relationships between f and τ, and between a and ¨q: τ = J Tf (7.6) and a = J ¨q + ˙J ˙q . (7.7) Combining these equations gives a = JH−1(J Tf −C) + ˙J ˙q = ΦAf + bA , (7.8) where ΦA = JH−1J T (7.9) and bA = ˙J ˙q −JH−1C . (7.10) These are the general equations for an articulated-body inverse inertia and bias acceleration. They apply to any articulated body, and they are independent of the choice of generalized coordinates. If the handle has a full six degrees of motion freedom, then ΦA is invertible, and we can define f = IAa + pA , (7.11) where IA = (JH−1J T)−1 (7.12) and pA = −IA bA . (7.13) Many of the basic properties of articulated-body inertias can be deduced from Eq. 7.12. For example, given that H is symmetric and positive-definite, it follows that IA must also be symmetric and positive-definite; and it is obvious that IA does not depend on velocity or force terms. Some basic properties of ΦA that follow from Eq. 7.9 are: rank(ΦA) = rank(J) , range(ΦA) = range(J) , (7.14) null(ΦA) = null(J T) .

A1 A2 A B1 B2 S fJ f1 a1 f2 a2 f a B1 Figure 7.4: Construction of a larger articulated body from two smaller ones Multiple Handles The projection method is easily extended to articulated bodies with more than one handle. Suppose there are h handles, and let ai, fi and Ji be the acceleration, force and Jacobian pertaining to handle i. If we redefine the quantities a, f and J to be a =

a1 ... ah

, f =

f1 ... fh

and J =

J1 ... Jh

, then Eqs. 7.6 to 7.14 immediately apply to the multiple-handle case. 7.2.2 The Assembly Method In the assembly method, we regard a given articulated body as an assembly of smaller articulated bodies; and we calculate its inertia and bias force from those of its component parts. Thus, the assembly method is suitable for use in recursive algorithms. In principle, the assembly method is general. However, we shall consider here only the special case in which every articulated body is a floating kinematic tree—a kinematic tree with no connection to a fixed base. This is the case that is needed for the articulated-body algorithm. The special properties of this case are: that every articulated body has one handle; that every handle has a full six degrees of freedom; and that every assembly operation consists of connecting one handle to another via a single joint. More general cases are covered in Sections 7.4 and 7.5. Figure 7.4 shows a pair of articulated bodies, A1 and A2, in the process of being assembled into a larger articulated body, A. The assembly operation consists of connecting the two handles, B1 and B2, with a joint. Prior to the assembly, the articulated-body equations for B1 and B2 are f1 = IA 1 a1 + pA 1 (7.15) and f2 = IA 2 a2 + pA 2 . (7.16)

After the assembly, B1 serves as the handle for the new articulated body, and its equation is f = IA a1 + pA . (7.17) The objective is to express IA and pA in terms of IA 1 , IA 2 , pA 1 , pA 2 and the motion constraint imposed by the new joint. When the assembly is made, the new joint introduces a new force, fJ, which is the total force transmitted through the joint from B1 to B2. The relationships between the force variables are therefore f1 = f −fJ , f2 = fJ . The value of fJ is unknown, but it imposes the motion constraint a2 −a1 = S¨q + c , and it satisfies STfJ = τ . In these equations, ¨q and τ are the joint’s acceleration and force variables, respectively, S is the joint’s motion subspace matrix, and c = ˙S ˙q (or ˙S ˙q + ˙σ if the joint depends on time). ¨q is unknown, but τ, S and c are given. By combining these equations with Eq. 7.16, we can calculate ¨q as follows: τ = STf2 = ST(IA 2 a2 + pA 2 ) = ST(IA 2 (a1 + c + S¨q) + pA 2 ) , so ¨q = (STIA 2 S)−1 (τ −ST(IA 2 (a1 + c) + pA 2 )) . (7.18) The matrix STIA 2 S is positive definite, and therefore invertible. Having solved for ¨q, we can now construct an equation relating f to a1 as follows: f = f1 + f2 = IA 1 a1 + IA 2 a2 + pA 1 + pA 2 = IA 1 a1 + IA 2 (a1 + c + S¨q) + pA 1 + pA 2 = IA 1 a1 + IA 2 (a1 + c + S (STIA 2 S)−1 (τ −ST(IA 2 (a1 + c) + pA 2 ))) + pA 1 + pA 2 . This equation has the same form as Eq. 7.17, and yields the following formulae for IA and pA: IA = IA 1 + IA 2 −IA 2 S (STIA 2 S)−1 STIA 2 (7.19)

and pA = pA 1 + pA 2 + IA 2 c + IA 2 S (STIA 2 S)−1 (τ −ST(IA 2 c + pA 2 )) . (7.20) With these formulae, we can compute the articulated-body inertia and bias force for any handle in any articulated body, provided only that the articulated body is a floating kinematic tree. Assembling Subtrees The three articulated bodies in Figure 7.4 were labelled A1, A2 and A merely for identification purposes, with no particular significance in the numbers 1 and 2. In practice, we would typically define a set of articulated bodies, A1 · · · ANB, such that Ai contained all of the bodies in the subtree rooted at body i, with body i serving as the handle (e.g. see Figure 7.5). In this circumstance, the most useful assembly formulae are the ones that tell how to calculate I A i and pA i from the inertias and bias forces pertaining to the children of body i; that is, the quantities IA j and pA j for all j ∈µ(i). Without repeating the algebra, these formulae are IA i = Ii + X j∈µ(i) Ia j (7.21) and pA i = pi + X j∈µ(i) pa j , (7.22) where Ii and pi are the rigid-body inertia and bias force for body i, and Ia j = IA j −IA j Sj (ST j IA j Sj)−1 ST j IA j (7.23) and pa j = pA j + Ia j cj + IA j Sj (ST j IA j Sj)−1 (τj −ST j pA j ) . (7.24) Essentially, these equations are obtained by applying Eqs. 7.19 and 7.20 once for each child of body i. On each application, A1 consists of body i plus every child subtree that has already been processed, and A2 is the next child subtree to be added. The quantities Ia j and pa j are the result of propagating IA j and pA j across joint j. They can be regarded as the apparent inertia and bias force of a massless handle that is connected to body j through joint j. In this case, the formulae in Eqs. 7.21 and 7.22 have the effect of gluing all of these massless handles to body i (see Figure 7.2). Ia j and pa j have one more interesting property: they express the force across a joint as a function of the parent body’s acceleration, whereas IA j and pA j express this same force as a function of the child’s acceleration: fj = IA j aj + pA j = Ia j aλ(j) + pa j . (7.25)

B1 B2 B3 = A3 A1 A2 S1 B1 f1 a1 (a) (b) Figure 7.5: Aspects of the articulated-body algorithm: (a) defining the articulated bodies, and (b) solving for the acceleration of joint 1 Ia j plays two roles in improving the efficiency of propagation calculations. First, it allows the two terms involving c in Eq. 7.20 to be collected into a single term, as shown in Eq. 7.24. Second, it has the property Ia j Sj = 0. Depending on the value of Sj, this property implies that one or more rows and columns of Ia j will be zero. For example, if Sj = [0 0 1 0 0 0]T (i.e., joint j is revolute), then row and column 3 of Ia j will be zero. This fact can be exploited to achieve a modest reduction in the cost of calculating articulated-body inertias. 7.3 The Articulated-Body Algorithm Suppose we have a kinematic tree containing NB rigid bodies. We can define a set of articulated bodies, A1 · · · ANB, such that Ai contains all of the bodies and joints in the subtree rooted at body i. A simple example is shown in Figure 7.5(a). Observe that body i (Bi) is in Ai but joint i is not. We define Bi to be the handle of Ai; and define IA i and pA i to be the articulated-body inertia and bias force for Bi in Ai. Consider the motion of B1, as shown in Figure 7.5(b). If f1 is the force transmitted across joint 1, then the equation of motion for B1 is f1 = IA 1 a1 + pA 1 . (7.26) Now, joint 1 constrains the acceleration of B1 to satisfy a1 = a0 + c1 + S1 ¨q1 , (7.27) where a0 is the known acceleration of the base, and c1 = ˙S1 ˙q1. The joint also constrains the transmitted force to satisfy ST 1 f1 = τ1 . (7.28)

From these three equations we get ST 1 (IA 1 (a0 + c1 + S1 ¨q1) + pA 1 ) = τ1 , which can be solved for ¨q1, giving ¨q1 = (ST 1 IA 1 S1)−1(τ1 −ST 1 IA 1 (a0 + c1) −ST 1 pA 1 )) . (7.29) Consider what we have just achieved. Simply knowing IA 1 and pA 1 has allowed us to calculate ¨q1 directly, without needing to know the values of any of the other acceleration variables. This is possible because Eq. 7.26 already accounts for the dynamic effect of every body in A1 on the acceleration response of B1. Having calculated ¨q1, it can be substituted into Eq. 7.27 to give the spatial acceleration of B1. This, in turn, allows us to repeat the process with the children of B1, and so on. In the general case, if the acceleration of body λ(i) is known, then the acceleration of body i and joint i can be calculated as follows: ¨qi = (ST i IA i Si)−1(τi −ST i IA i (aλ(i) + ci) −ST i pA i )) (7.30) and ai = aλ(i) + ci + Si ¨qi . (7.31) Therefore, if the articulated-body inertias and bias forces are known, then the forward dynamics of the kinematic tree can be calculated using these two equations, with i ranging from 1 to NB. Algorithm The articulated-body algorithm calculates the forward dynamics of a kinematic tree in O(NB) arithmetic operations, in exactly the manner outlined above. It makes three passes over the tree, as follows: an outward pass (root to leaves) to calculate velocity and bias terms; an inward pass to calculate articulated-body inertias and bias forces; and a second outward pass to calculate the accelerations. Pass 1 The job of the first pass is to calculate the velocity-product accelerations, ci, and the rigid-body bias forces, pi, for use in later passes. Both are functions of the body velocities, vi, so it is necessary to calculate these as well. The equations of the first pass are: vJi = Si ˙qi , (7.32) cJi = ˚ Si ˙qi , (7.33) vi = vλ(i) + vJi , (v0 = 0) (7.34) ci = cJi + vi × vJi , (7.35) pi = vi ×∗Ii vi −f x i . (7.36)

They are similar to the equations of the first pass in the recursive Newton-Euler algorithm, as shown in Table 5.1. The vectors vJi and cJi will have values of Si ˙qi + σi and ˚ Si ˙qi + ˚σi, respectively, if joint i depends explicitly on time (see §3.5); but the calculation details are handled by jcalc (§4.4). f x i is the external force, if any, acting on body i. Pass 2 The next step is to calculate the articulated-body inertias and bias forces, IA i and pA i . This is done using the assembly formulae in Eqs. 7.21 to 7.24, which are reproduced below. IA i = Ii + P j∈µ(i) Ia j (7.37) pA i = pi + P j∈µ(i) pa j , (7.38) Ia j = IA j −IA j Sj(ST j IA j Sj)−1ST j IA j (7.39) pa j = pA j + Ia j cj + IA j Sj(ST j IA j Sj)−1(τj −ST j pA j ) . (7.40) These calculations are performed for each value of i from NB down to 1. This implies that Ia j and pa j are not calculated for every j, but only for those values that satisfy λ(j)̸ = 0. If body i has no children then IA i = Ii and pA i = pi. Pass 3 The third pass calculates the accelerations using Eqs. 7.30 and 7.31, which we repeat here as ¨qi = (ST i IA i Si)−1(τi −ST i IA i (aλ(i) + ci) −ST i pA i )) (7.41) ai = aλ(i) + ci + Si ¨qi . (a0 = −ag) (7.42) By initializing the base acceleration to −ag, rather than 0, we can simulate the effect of a uniform gravitational field with a fictitious upward acceleration. This is more efficient than treating gravity as an external force. Common Subexpressions Equations 7.39 to 7.42 contain a number of common subexpressions, which can be eliminated by defining the following intermediate results: Ui = IA i Si , (7.43) Di = ST i Ui , (7.44) ui = τi −ST i pA i , (7.45) a′ i = aλ(i) + ci . (7.46) In terms of these quantities, Eqs. 7.39 to 7.42 simplify to Ia j = IA j −Uj D−1 j U T j , (7.47) pa j = pA j + Ia j cj + Uj D−1 j uj , (7.48) ¨qi = D−1 i (ui −U T i a′ i) , (7.49) ai = a′ i + Si ¨qi . (a0 = −ag) (7.50)

Equations in Body Coordinates The articulated-body algorithm works best in body coordinates. To convert the above equations to body coordinates, simply define each quantity that pertains to body i as being expressed in body i coordinates, and insert coordinate transforms where they are needed. The resulting equations are shown in Table 7.1. Algorithm Details The pseudocode for the articulated-body algorithm is shown in Table 7.1. It is organized into three passes, with one for loop for each pass. The quantities XJ, vJ and cJ are local variables in the first pass; Ia and pa are local variables in the second pass; and a′ is a local variable in the third pass. If the tree contains joints that depend explicitly on time, then add time as a fourth argument to jcalc. If the tree contains joints for which the velocity variables are not the derivatives of the position variables, then replace the symbols ˙qi and ¨qi with αi and ˙αi. If there are no external forces, then the calculation of iX0 is not necessary. The pseudocode assumes that as soon as a transform X has been calculated, the transforms X∗, X−1 and (X∗)−1 are immediately available for use. The functions that make this possible are described in Appendix A. Equations 7.37 and 7.38 have been transformed from calculations using µ(i) to calculations using λ(i). As a result, IA i and pA i are initialized to Ii and pi at the bottom of the first for loop, and the remainder of the calculation is performed inside the if statement in the second for loop (cf. the implementation of Eq. 5.20 in §5.3). 7.4 Alternative Assembly Formulae Section 7.3 describes the standard version of the articulated-body algorithm, which is based on the assembly formulae in Section 7.2.2. However, there are other formulae that could be used instead. This section presents a systematic procedure for deriving alternative assembly formulae, which give rise to alternative versions of the articulated-body algorithm. In most cases, the alternatives are more complicated and less efficient than the standard algorithm. However, they can solve a larger class of problems. In Section 7.2.2 we solved the following problem: given the five equations f1 = IA 1 a1 + pA 1 , (7.51) f2 = IA 2 a2 + pA 2 , (7.52) f = f1 + f2 , (7.53) a2 = a1 + c + S ¨q (7.54) and STf2 = τ , (7.55)

Equations (in body coordinates): Pass 1 v0 = 0 vJi = Si ˙qi cJi = ˚ Si ˙qi vi = iXλ(i) vλ(i) + vJi ci = cJi + vi × vJi pi = vi ×∗Ii vi −iX∗ 0 f x i Pass 2 IA i = Ii + X j∈µ(i) iX∗ j Ia j jXi pA i = pi + X j∈µ(i) iX∗ j pa j Ui = IA i Si Di = ST i Ui ui = τi −ST i pA i Ia i = IA i −Ui D−1 i U T i pa i = pA i + Ia i ci + Ui D−1 i ui Pass 3 a0 = −ag a′ i = iXλ(i) aλ(i) + ci ¨qi = D−1 i (ui −U T i a′ i) ai = a′ i + Si ¨qi Algorithm: v0 = 0 for i = 1 to NB do [XJ, Si, vJ, cJ] = jcalc(jtype(i), qi, ˙qi) iXλ(i) = XJ XT(i) if λ(i)̸ = 0 then iX0 = iXλ(i) λ(i)X0 end vi = iXλ(i) vλ(i) + vJ ci = cJ + vi × vJ IA i = Ii pA i = vi ×∗Ii vi −iX∗ 0 f x i end for i = NB to 1 do Ui = IA i Si Di = ST i Ui ui = τi −ST i pA i if λ(i)̸ = 0 then Ia = IA i −Ui D−1 i U T i pa = pA i + Iaci + Ui D−1 i ui IA λ(i) = IA λ(i) + λ(i)X∗ i Ia iXλ(i) pA λ(i) = pA λ(i) + λ(i)X∗ i pa end end a0 = −ag for i = 1 to NB do a′ = iXλ(i) aλ(i) + ci ¨qi = D−1 i (ui −U T i a′) ai = a′ + Si ¨qi end Table 7.1: The articulated body equations and algorithm

find expressions for the coefficients IA and pA in the target equation f = IA a1 + pA . (7.56) By following a few simple rules, the five given equations can be assembled into the following composite equation:

1 −1 −1 0 0 0 1 0 0 0 0 0 1 −IA 2 0 0 0 0 1 −S 0 0 ST 0 0

f f1 f2 a2 ¨q

=

0 pA 1 pA 2 c τ

+

0 IA 1 0 1 0

a1 . The rules for constructing this equation are: 1. the right-hand unknown in the target equation (a1) becomes right-hand unknown in the composite equation; 2. the left-hand unknown in the target equation (f) becomes the top entry in the composite vector of unknowns; 3. the unknown vector with a dimension other than six (¨q) becomes the bottom entry in the composite vector of unknowns; 4. the given equation with a dimension other than six (STf2 = τ) occupies the bottom row in the composite equation; and 5. the remaining equations and unknowns are inserted in any order that results in the 5 × 5 composite matrix being as close as possible to uppertriangular form. Having constructed this equation, the original problem can be solved using a process of Gaussian elimination and back-substitution. Once the coefficient matrix has been brought into upper-triangular form, we immediately have an equation expressing ¨q as a function of a1; and the process of back-substitution eventually yields an equation giving f as a function of a1. At this point, the coefficient of a1 is the desired expression for IA, and the remaining terms provide the desired expression for pA. The process is illustrated in Table 7.2 under the heading ‘Standard Formulae’. At the end of the Gaussian elimination stage, the fifth row reads STIA 2 S¨q = τ −STpA 2 −STIA 2 c −STIA 2 a1 , which implies ¨q = (STIA 2 S)−1(τ −STpA 2 −STIA 2 c −STIA 2 a1) (cf. Eq. 7.18). Having found ¨q in terms of a1, back-substitution allows us to express a2 in terms of a1, and so on, until eventually we get an expression for f in terms of a1, from which the formulae for IA and pA can be extracted.

Standard Formulae obtain f = IAa1 + pA using IA i , pA i and S 2 66664 1 −1 −1 0 0 0 1 0 0 0 0 0 1 −IA 2 0 0 0 0 1 −S 0 0 ST 0 0 3 77775 2 66664 f f1 f2 a2 ¨q 3 77775 = 2 66664 0 pA 1 pA 2 c τ 3 77775 + 2 66664 0 IA 1 0 1 0 3 77775 a1 2 66664 1 −1 −1 0 0 0 1 0 0 0 0 0 1 −IA 2 0 0 0 0 1 −S 0 0 0 STIA 2 0 3 77775 2 66664 f f1 f2 a2 ¨q 3 77775 = 2 66664 0 pA 1 pA 2 c τ −STpA 2 3 77775 + 2 66664 0 IA 1 0 1 0 3 77775 a1 2 66664 1 −1 −1 0 0 0 1 0 0 0 0 0 1 −IA 2 0 0 0 0 1 −S 0 0 0 0 STIA 2 S 3 77775 2 66664 f f1 f2 a2 ¨q 3 77775 = 2 66664 0 pA 1 pA 2 c z 3 77775 + 2 66664 0 IA 1 0 1 −STIA 2 3 77775 a1 where z = τ −STpA 2 −STIA 2 c result: ¨q = (STIA 2 S)−1(z −STIA 2 a1) IA = IA 1 + IA 2 −IA 2 S(STIA 2 S)−1STIA 2 pA = pA 1 + pA 2 + IA 2 c + IA 2 S(STIA 2 S)−1z Variant Formulae 1 obtain a1 = ΦAf + bA using ΦA i , bA i and T 2 66664 1 0 −ΦA 1 0 0 0 1 0 −ΦA 2 0 0 0 1 1 0 0 0 0 1 −T −T T T T 0 0 0 3 77775 2 66664 a1 a2 f1 f2 λ 3 77775 = 2 66664 bA 1 bA 2 0 Taτ T Tc 3 77775 + 2 66664 0 0 1 0 0 3 77775 f 2 66664 1 0 −ΦA 1 0 0 0 1 0 −ΦA 2 0 0 0 1 1 0 0 0 0 1 −T 0 0 −T TΦA 1 T TΦA 2 0 3 77775 2 66664 a1 a2 f1 f2 λ 3 77775 = 2 66664 bA 1 bA 2 0 Taτ T T(c + bA 1 −bA 2 ) 3 77775 + 2 66664 0 0 1 0 0 3 77775 f (continued on next page) Table 7.2: Alternative assembly formulae

(continued from previous page) 2 66664 1 0 −ΦA 1 0 0 0 1 0 −ΦA 2 0 0 0 1 1 0 0 0 0 1 −T 0 0 0 T T(ΦA 1 + ΦA 2 ) 0 3 77775 2 66664 a1 a2 f1 f2 λ 3 77775 = 2 66664 bA 1 bA 2 0 Taτ T T(c + bA 1 −bA 2 ) 3 77775 + 2 66664 0 0 1 0 T TΦA 1 3 77775 f 2 66664 1 0 −ΦA 1 0 0 0 1 0 −ΦA 2 0 0 0 1 1 0 0 0 0 1 −T 0 0 0 0 T T(ΦA 1 + ΦA 2 )T 3 77775 2 66664 a1 a2 f1 f2 λ 3 77775 = 2 66664 bA 1 bA 2 0 Taτ z 3 77775 + 2 66664 0 0 1 0 T TΦA 1 3 77775 f where z = T T(c + bA 1 −bA 2 −(ΦA 1 + ΦA 2 )Taτ) result: λ = (T T(ΦA 1 + ΦA 2 )T )−1(z + T TΦA 1 f) ΦA = ΦA 1 −ΦA 1 T (T T(ΦA 1 + ΦA 2 )T )−1T TΦA 1 bA = bA 1 −ΦA 1 Taτ −ΦA 1 T (T T(ΦA 1 + ΦA 2 )T )−1z Variant Formulae 2 obtain f = IAa1 + pA using IA i , pA i and T 2 66664 1 −1 −1 0 0 0 1 0 0 0 0 0 1 −IA 2 0 0 0 1 0 −T 0 0 0 T T 0 3 77775 2 66664 f f1 f2 a2 λ 3 77775 = 2 66664 0 pA 1 pA 2 Taτ T Tc 3 77775 + 2 66664 0 IA 1 0 0 T T 3 77775 a1 2 66664 1 −1 −1 0 0 0 1 0 0 0 0 0 1 −IA 2 0 0 0 0 IA 2 −T 0 0 0 T T 0 3 77775 2 66664 f f1 f2 a2 λ 3 77775 = 2 66664 0 pA 1 pA 2 Taτ −pA 2 T Tc 3 77775 + 2 66664 0 IA 1 0 0 T T 3 77775 a1 2 66664 1 −1 −1 0 0 0 1 0 0 0 0 0 1 −IA 2 0 0 0 0 IA 2 −T 0 0 0 0 T T(IA 2 )−1T 3 77775 2 66664 f f1 f2 a2 λ 3 77775 = 2 66664 0 pA 1 pA 2 Taτ −pA 2 z 3 77775 + 2 66664 0 IA 1 0 0 T T 3 77775 a1 where z = T T(c −(IA 2 )−1(Taτ −pA 2 )) result: λ = (T T(IA 2 )−1T )−1(z + T Ta1) IA = IA 1 + T (T T(IA 2 )−1T )−1T T pA = pA 1 + Taτ + T (T T(IA 2 )−1T )−1z Table 7.2: Alternative assembly formulae

The point of this exercise is that it makes the solution process systematic, so that it can be applied to a different set of given equations. There are two obvious changes we can try: 1. replace Eqs. 7.51, 7.52 and 7.56 with their inverses; that is, a1 = ΦA 1 f1 + bA 1 , a2 = ΦA 2 f2 + bA 2 and a1 = ΦAf + bA; and 2. replace Eqs. 7.54 and 7.55 with the equations of the equivalent implicit joint constraint, which are T Ta2 = T T(a1 + c) and f2 = Taτ + T λ. Table 7.2 shows two of the alternative formulae that can be produced by this method: variant 1 is the result of making both changes; and variant 2 is the result of making only the second change. Each can be used as the basis for an alternative formulation of the articulated-body algorithm.1 Observe that this process gives rise to some potentially useful matrix identities. For example, the variant formula for IA must be equal to the standard formula, and both must equal the inverse of the formula for ΦA. Similar comments apply to pA and bA. Any formula that uses inertias, rather than inverse inertias, is restricted by the requirement that the inertias must exist. Thus, the standard formulae, and those of variant 2, are only applicable if both handles have a full six degrees of freedom. In contrast, the formulae of variant 1 require only that the matrix T T(ΦA 1 + ΦA 2 )T have full rank. Physically, this is equivalent to requiring that the joint does not impose a redundant constraint; that is, it does not duplicate a constraint that was already there. A sufficient condition for the application of variant 1 is that at least one of the two handles have a full six degrees of freedom. Thus, an algorithm based on variant 1 would be able to calculate articulatedbody inverse inertias for bodies in any kinematic tree, and a few closed-loop systems, whereas the standard algorithm is limited to floating kinematic trees. 7.5 Multiple Handles So far, we have applied the assembly method only to articulated bodies having one handle each. This is sufficient for assembling kinematic trees. However, if we wish to assemble general rigid-body systems, or if we want greater flexibility in the order of assembly, then we must be able to assemble articulated bodies having more than one handle. As an example of what can be accomplished, Featherstone (1999a,b) describes a divide-and-conquer algorithm to assemble both kinematic trees and closed-loop systems in a manner that allows their dynamics to be calculated in O(log(NB)) time on a parallel computer with O(NB) processors. Therefore, let us now develop an assembly formula for articulated bodies with multiple handles. 1Neither variant calculates ¨q as a by-product; but this quantity can be calculated from a1 and a2 using the formula ¨q = T T a (a2 −a1 −c). An algorithm based on variant 1 will be more efficient if the active joint force, Taτ, is replaced by an equivalent pair of external forces which contribute to the bias accelerations of the individual rigid bodies.

A12 A34 A S fJ f1 a1 f4 a4 1 2 3 4 Figure 7.6: Assembly of articulated bodies with multiple handles Figure 7.6 shows two articulated bodies, called A12 and A34, which are to be assembled to form a larger articulated body called A. Body A12 contains two handles, numbered 1 and 2, while A34 contains handles 3 and 4. The idea is to assemble these bodies by connecting a joint between handles 2 and 3. The assembly operation is deemed to consume these two handles, leaving articulated body A with only handles 1 and 4. We shall assume that at least one of handles 2 and 3 has six degrees of freedom. The equations of motion for A12 and A34 before the assembly operation are  a1 a2  =  Φ1 Φ12 Φ21 Φ2   f1 f2  +  b1 b2  (7.57) and a3 a4  =  Φ3 Φ34 Φ43 Φ4  f3 f4  + b3 b4  ; (7.58) and the equation of motion for A after the assembly operation is a1 a4  = ΦA 1 ΦA 14 ΦA 41 ΦA 4  f1 f4  + bA 1 bA 4  . (7.59) (Observe that we have simplified the notation by omitting the superscripts from the coefficients of Eqs. 7.57 and 7.58.) We are using inverse inertias here because it is possible—and indeed quite likely—that the two handles within an articulated body do not have a full 12 degrees of motion freedom between them, even if they do have 6 degrees of freedom individually. In the absence of full motion freedom, the coefficient matrices in these equations will be singular. However, our assumption that at least one of the two handles 2 and 3 has six degrees of freedom on its own implies that at least one of the two matrices Φ2 and Φ3 is positive definite. This, in turn, implies that the sum Φ2 + Φ3 is positive definite.

When the assembly is made, the joint will impose the following motion constraint on handles 2 and 3: T T(a3 −a2) = T Tc , (7.60) where c = ˙S ˙q, and it will also impose the force constraint f3 = T λ = −f2 . (7.61) In formulating Eq. 7.61, we have assumed that any active force at the joint has already been accounted for in the bias accelerations of Eqs. 7.57 and 7.58. In other words, the active joint force, Taτ, has been replaced by a pair of equal and opposite external forces, Taτ acting on handle 3 and −Taτ acting on handle 2, and that the effects of these two forces on the motion of A12 and A34 have already been included in the vectors b1 · · · b4. The objective is to express the coefficients of Eq. 7.59 in terms of the coefficients of Eqs. 7.57, 7.58, 7.60 and 7.61. The first step in the solution process is to substitute for a3 and a2 in Eq. 7.60, giving T T(Φ3f3 + Φ34f4 + b3 −Φ21f1 −Φ2f2 −b2) = T Tc . The next step is to substitute for f3 and f2 in this equation, using Eq. 7.61, and collect terms in λ, which gives T T(Φ2 + Φ3)T λ = T T(Φ21f1 −Φ34f4 + c −b3 + b2) . Given that Φ2 + Φ3 is positive definite, the coefficient of λ is guaranteed to be nonsingular, and so we may solve for λ as follows: λ = Y −1T T(Φ21f1 −Φ34f4 + β) , (7.62) where Y = T T(Φ2 + Φ3)T and β = c −b3 + b2 . The final step is to substitute for f2 and f3 in the expressions for a1 and a4 (drawn from Eqs. 7.57 and 7.58), using both Eq. 7.61 and Eq. 7.62. This gives a1 = Φ1f1 −Φ12T Y −1T T(Φ21f1 −Φ34f4 + β) + b1 a4 = Φ4f4 + Φ43T Y −1T T(Φ21f1 −Φ34f4 + β) + b4 . These two equations express a1 and a4 as functions of f1 and f4 only, and are therefore the component parts of Eq. 7.59. Thus, the formulae for the coefficients of Eq. 7.59 are ΦA 1 = Φ1 −Φ12W Φ21 ΦA 4 = Φ4 −Φ43W Φ34 ΦA 14 = Φ12W Φ34 = (ΦA 41)T bA 1 = b1 −Φ12W β bA 4 = b4 + Φ43W β (7.63) where W = T Y −1T T.

Generalizations To generalize the number of handles, we replace the four individual handles with four handle sets. Set 1 contains every handle in A12 that will not be connected to a joint during assembly; set 2 contains every handle that will; and so on. Each handle in set 2 will be connected to one handle in set 3 via one joint; so sets 2 and 3 must contain the same number of handles. To permit multiple joints to be connected to the same body, we allow a body to have more than one handle. In this context, a handle should no longer be regarded as a body, but as a location on a body. Any kinematic tree can be obtained by a sequence of assembly operations involving one joint per assembly (i.e., sets 2 and 3 contain one handle each). However, the construction of closed-loop systems requires multi-joint assemblies. Having defined these sets, the next step is to define a1, f2, Φ34, and so on, to be composite vectors and matrices containing every relevant vector and matrix for their sets. For example, if set 1 contains handles 1a, 1b, 1c, . . . , having accelerations of a1a, a1b, . . . , then a1 is the composite vector [aT 1a aT 1b · · · ]T. Note that T will be a block-diagonal matrix composed of the constraint-force subspace matrices of the individual joints. Given these composite vectors and matrices, the derivation of Eq. 7.63 proceeds in the same manner as shown in Eqs. 7.57 to 7.62, except that it may no longer be possible to assume that Y is invertible. Featherstone (1999b) explains what to do if Y is singular. It also explains how to incorporate constraint stabilization terms (see §8.3) into the assembly formula.
